\section{Заключение}

В ходе выполнения лабораторной работы была смоделирована схема источника электрической энергии в среде LTspice и проведено исследование его характеристик при различных значениях сопротивления нагрузки.

По результатам обработки экспериментальных данных и теоретических расчётов были сделаны следующие выводы касательно режимов работы источника электрической энергии постоянного тока:

\begin{enumerate}
    \item \textit{Режим холостого хода:} ток в цепи равен нулю, напряжение на зажимах достигает максимального значения и равно ЭДС, а коэффициент полезного действия принимает максимальное значение.
    \item \textit{Режим короткого замыкания:} напряжение на зажимах падает до $0~\text{В}$, ток в цепи равен предельному значению $J_{\mathrm{p}}$, вся мощность рассеивается на внутреннем сопротивлении источника, а КПД равен $0\%$.
    \item \textit{Согласованный режим:} в нагрузке выделяется максимальная полезная мощность $P_{\max}$. При этом напряжение на нагрузке составляет половину ЭДС, а КПД источника равен $50\%$.
\end{enumerate}

Построенные внешняя $U(I)$ и рабочие $P(I)$, $\eta(I)$ характеристики продемонстрировали сходимость результатов моделирования в LTspice с аналитическими расчётами.